% !TEX TS-program = luatex
% Jonas Heinle -- CV
%
% Layout: myCV_METADATA class (md2pdfLib/cv/template/latex/), found via
% TEXINPUTS. Language (english/german) is selected by the build, see
% data/cv/README.md.
%
% Based on a template created by Sabrina Benge, itself derived from
% Christophe Roger's YAAC / Awesome Source CV.
% Template license: CC BY-SA 4.0 (https://creativecommons.org/licenses/by-sa/4.0/)

\documentclass[alternative]{myCV_METADATA}

% Brand CV environments: timeline, language proficiency.
% Found via TEXINPUTS (see scripts/build_in_container.sh).
\tcbuselibrary{breakable}
\input{brand-cv.tex}

% Profile: which sections this build shows and the tagline above them. One CV,
% several targets -- the general-purpose one published on jonasheinle.de, and
% per-application variants that reorder the same sections and swap the summary.
% Selected at build time like the language, without editing this file:
%   lualatex "\def\cvprofile{mistral-rse}% !TEX TS-program = luatex
% Jonas Heinle -- CV
%
% Layout: myCV_METADATA class (md2pdfLib/cv/template/latex/), found via
% TEXINPUTS. Language (english/german) is selected by the build, see
% data/cv/README.md.
%
% Based on a template created by Sabrina Benge, itself derived from
% Christophe Roger's YAAC / Awesome Source CV.
% Template license: CC BY-SA 4.0 (https://creativecommons.org/licenses/by-sa/4.0/)

\documentclass[alternative]{myCV_METADATA}

% Brand CV environments: timeline, language proficiency.
% Found via TEXINPUTS (see scripts/build_in_container.sh).
\tcbuselibrary{breakable}
\input{brand-cv.tex}

% Profile: which sections this build shows and the tagline above them. One CV,
% several targets -- the general-purpose one published on jonasheinle.de, and
% per-application variants that reorder the same sections and swap the summary.
% Selected at build time like the language, without editing this file:
%   lualatex "\def\cvprofile{mistral-rse}% !TEX TS-program = luatex
% Jonas Heinle -- CV
%
% Layout: myCV_METADATA class (md2pdfLib/cv/template/latex/), found via
% TEXINPUTS. Language (english/german) is selected by the build, see
% data/cv/README.md.
%
% Based on a template created by Sabrina Benge, itself derived from
% Christophe Roger's YAAC / Awesome Source CV.
% Template license: CC BY-SA 4.0 (https://creativecommons.org/licenses/by-sa/4.0/)

\documentclass[alternative]{myCV_METADATA}

% Brand CV environments: timeline, language proficiency.
% Found via TEXINPUTS (see scripts/build_in_container.sh).
\tcbuselibrary{breakable}
\input{brand-cv.tex}

% Profile: which sections this build shows and the tagline above them. One CV,
% several targets -- the general-purpose one published on jonasheinle.de, and
% per-application variants that reorder the same sections and swap the summary.
% Selected at build time like the language, without editing this file:
%   lualatex "\def\cvprofile{mistral-rse}% !TEX TS-program = luatex
% Jonas Heinle -- CV
%
% Layout: myCV_METADATA class (md2pdfLib/cv/template/latex/), found via
% TEXINPUTS. Language (english/german) is selected by the build, see
% data/cv/README.md.
%
% Based on a template created by Sabrina Benge, itself derived from
% Christophe Roger's YAAC / Awesome Source CV.
% Template license: CC BY-SA 4.0 (https://creativecommons.org/licenses/by-sa/4.0/)

\documentclass[alternative]{myCV_METADATA}

% Brand CV environments: timeline, language proficiency.
% Found via TEXINPUTS (see scripts/build_in_container.sh).
\tcbuselibrary{breakable}
\input{brand-cv.tex}

% Profile: which sections this build shows and the tagline above them. One CV,
% several targets -- the general-purpose one published on jonasheinle.de, and
% per-application variants that reorder the same sections and swap the summary.
% Selected at build time like the language, without editing this file:
%   lualatex "\def\cvprofile{mistral-rse}\input{cv.tex}"
% which is what CV_PROFILE in scripts/build_in_container.sh does. Every profile
% defines \cvTagline and \cvBody; see profiles/README.md.
%
% Profiles select *sections*, they do not fork prose: a per-profile branch
% inside every section file is the same trap data/cv/README.md flags for a
% third language. Role-specific wording lives in its own section_headline_*
% file, which the profile inputs instead of the general one.
%
% The one exception to "profiles do not edit sections": the A-levels entry in
% section_education.tex. A profile aimed at a role that requires a Master's
% spends three lines on a 2006 Gymnasium entry it does not need, and the page
% budget is the whole constraint. A \newif rather than an \ifdefstring: the
% entry it guards contains blank lines, which an argument-grabbing conditional
% would choke on. Declared here, before the profile is read, so a profile can
% flip it.
\newif\ifcvshowschool \cvshowschooltrue
% Same idea, one entry up: drops the thesis lines while keeping the degrees,
% their dates, their grades and their subject tags.
\newif\ifcvshowthesislines \cvshowthesislinestrue
% Suppresses the "Teaching and Mentoring" heading so the tutor entry can follow
% the other jobs under Experiences -- which is what it is, and which is what
% makes an unbroken 2018-today employment line visible on one page. Only useful
% for a profile that inputs section_teaching_mentoring directly after
% section_experience.
\newif\ifcvshowteachingtitle \cvshowteachingtitletrue
% Drops the CUDA-kernels bullet. The one bullet a platform or infrastructure
% posting does not buy, and the cheapest two lines in the Fraunhofer entry for
% a profile that needs room. On for every research-facing profile.
\newif\ifcvshowcudakernels \cvshowcudakernelstrue
% The AI-gateway bullet, and the only flag here that defaults OFF: it is the
% one place the CV claims Go, so a profile opts in rather than every existing
% profile paying two lines it did not budget for. \ifcvrelocate is the same
% shape. `all` turns it on, which is what `all` means.
\newif\ifcvshowgateway \cvshowgatewayfalse
% Drops the RemaNet bullet -- three lines, the longest in the entry. Edge-AI
% product work: central to a research profile, the first thing a platform
% posting trades away.
\newif\ifcvshowremanet \cvshowremanettrue
% The documentation bullet. Default off for the same reason as the gateway:
% new content added for one application, and every other profile is at zero
% slack. `all` turns it on.
\newif\ifcvshowdocs \cvshowdocsfalse
% The quantum-physics GPU-port bullet (Fraunhofer IPA's quantum-computing
% group). Same shape as the gateway/docs pair: new content for one application
% -- the AMD HPC one -- and the published profiles are at zero slack. `all`
% turns it on. The amd-hpc profile trades the RemaNet bullet for this one.
\newif\ifcvshowquantum \cvshowquantumfalse
% The SemProbe one-liner under the ECML PKDD entry in section_private.tex. It
% used to be commented out and is now behind this flag. Default off because it
% costs two \footnotesize lines in Publications, the column that sets the
% height of the bottom row, and the published profiles have zero slack (the
% German mistral-rse build is the one it pushes over). terabase-cv opts in
% because it is the only section line that says "object detectors", and that
% profile's summary sentence on detectors rests on it. `all` turns it on.
\newif\ifcvshowsemprobe \cvshowsemprobefalse
% The Berghof bullet: object detection on an embedded Berghof PLC fitted with
% a Hailo accelerator. Until 2026-09-22 it sat inside the RemaNet bullet, which
% made it read as part of RemaNet. The owner corrected that: it was a separate
% project, with its own full ML lifecycle. Default ON, so every profile that
% showed the PLC still shows it. amd-hpc turns it off: that profile hides
% RemaNet, so it never showed the PLC.
\newif\ifcvshowberghof \cvshowberghoftrue
% The industrial safety project: person detection, 4 cameras at once, their
% inferences batched, full ML lifecycle (owner, 2026-09-22). New content, so
% it is default OFF like gateway/docs/quantum: the published profiles have
% zero slack. `all` and terabase-cv turn it on.
\newif\ifcvshowsafety \cvshowsafetyfalse
% The build-systems bullet: CMake, GStreamer (Meson), LiteRT-LM (Bazel), and
% fixes upstreamed to LLVM and sccache. On everywhere except terabase-cv,
% where its three lines pay for the Berghof and safety bullets. It is the
% bullet a computer-vision posting has the least use for.
\newif\ifcvshowbuilds \cvshowbuildstrue
% The KI-PrOT bullet: visual classification of parts into drag-out classes
% (surface treatment; ki-prot.de, Fraunhofer IPA develops the AI methods).
% The owner built and annotated the dataset and versioned it with DVC
% (owner, 2026-09-22). New content, so default OFF; terabase-cv and `all`
% turn it on.
\newif\ifcvshowkiprot \cvshowkiprotfalse
% ML-lifecycle detail clauses, appended to existing bullets when on: INT8 and
% the customer field-feedback loop on the Berghof PLC, the owner's own
% self-labeling research tool (some of the owner's projects used its generated
% labels -- reworded 2026-09-23; first put as "its labels partly trained the
% models"), and
% MLflow/ClearML experiment tracking. All owner-confirmed on 2026-09-22. Default
% OFF: the base sentences stay exactly as they were, so zero-slack profiles
% are untouched. terabase-cv and `all` turn it on.
\newif\ifcvshowmldetail \cvshowmldetailfalse
% Stack-detail clauses, appended when on (owner, 2026-09-23): the coding
% agent's open-weight models come from Hugging Face, and RemaNet's web part
% used TypeScript and Node.js. Since the same day it also names Docker in the
% CI/CD bullet ("cut Docker builds"): Docker daily at Fraunhofer, advanced
% builds (owner). Default OFF, and when off the base sentences
% typeset exactly as before, so no other profile's lines move. parallel-ai
% and `all` turn it on.
\newif\ifcvshowstackdetail \cvshowstackdetailfalse
% Drops the CAS description, keeping the entry: its dates, title, employer and
% tags all stay, so the timeline is untouched. The ONNX/mixed-reality bullet is
% the one description on the page a platform posting has no use for. Guards the
% whole itemize, not the \item -- an itemize with no \item is a LaTeX error.
\newif\ifcvshowcasdetail \cvshowcasdetailtrue
% Off by default, on for a profile aimed at a job in another city. The header
% is where a reader checks whether the geography works, and it answers the
% question "Stuttgart, but the job is in Paris?" in the same place it is asked.
% Costs nothing: with the birthday and the YouTube handle gone both contact
% columns run four lines, and the citation column beside them is five (English)
% or seven (German), so it is the citation that sets the header height. Only
% grow this past five lines with the German build in front of you.
\newif\ifcvrelocate \cvrelocatefalse
% The order of the Fraunhofer IPA bullets. section_experience.tex defines each
% one as \cvFhg<Name> (each still behind its own \ifcvshow... flag) and prints
% them through this macro, so a profile can lead with what its posting asks
% for first without forking any prose: \renewcommand it in the profile.
% This default is the order the section file has always had. List every
% bullet in a profile's order too, even hidden ones: a bullet left out of the
% list never prints, whatever its flag says, and that failure is silent.
\newcommand{\cvFraunhoferOrder}{%
	\cvFhgKernels\cvFhgQuantum\cvFhgCicd\cvFhgBuilds\cvFhgGateway\cvFhgDocs%
	\cvFhgVlm\cvFhgAgent\cvFhgRemanet\cvFhgBerghof\cvFhgSafety\cvFhgKiprot}
% The Fraunhofer IPA tag row, set per profile the same way (2026-09-23). The
% row is one line wide and every tag in it competes for that line, so a
% profile whose posting has no use for CMake, Bazel and Meson -- the build
% systems of a bullet it hides -- can spend the room on what it does ask for:
% \renewcommand it in the profile. This default is the list the section file
% had before, so no other profile changes. The \experience override below
% expands the tag argument once, which is what lets a macro stand in for a
% literal list there.
\newcommand{\cvFhgTags}{Python, uv, Go, C++, CUDA, PyTorch, CMake, Bazel, Meson, vLLM, LoRA, GitLab CI/CD}
% \cvshowthesislinesfalse hid the thesis text but not the row holding it:
% \education always emits its description row, so switching the flag off
% left an empty minipage -- a blank line under every degree, which is exactly
% the vertical space the flag was set to reclaim. Emitting that row only when
% there is something to put in it makes the flag pay what it promises.
% Overridden here rather than in myCV_METADATA.cls, which other documents
% share, following the \socialtext/\sociallink precedent below.
\renewcommand\education[8]{
	\textbf{#1}    & \textbf{#2, \textsc{#3}, #4}   \\*
	\textbf{#5}    & #6   \\*
	\ifcvshowthesislines
	& \begin{minipage}[t]{\rightcolumnlength}
		#7
	\end{minipage}   \\*
	\fi
	& \footnotesize{\foreach \n in {#8}{\cvtag{\n}}} \\
}
% Extra space between an experience's description and its tag row. The class
% sets the tag boxes directly under the last description line, and when that
% line has descenders the boxes nearly touch it ("much too less space", owner,
% 2026-09-22, on the Fraunhofer entry). \cvtagsep is 0pt by default, and at
% zero the row ends in the class's plain \\* -- token for token the class's own
% definition, so no profile changes unless it sets the length. terabase-cv
% does. \expandafter drops the \else/\fi before \\ runs, so no conditional is
% left open across the row end.
% The class's \experience, copied verbatim below except for the one row end.
% Overridden here, not in myCV_METADATA.cls, for the same reason as
% \education above.
\newlength{\cvtagsep}
\setlength{\cvtagsep}{0pt}
\makeatletter
% How it works (third attempt, 2026-09-22): a zero-width rule of height
% \cvtagsep closes the description minipage. It deepens that cell by exactly
% \cvtagsep, the timeline's vertical line stays unbroken (same row), and an
% empty description (a hidden CAS detail) gains nothing: there the rule is
% the minipage's first item and only sets its height, which the row strut
% already exceeds. At 0pt it expands to nothing, so the class layout is
% unchanged token for token.
\newcommand\cv@descpad{\ifdim\cvtagsep>\z@
		\par\nointerlineskip\hrule\@height\cvtagsep\@width\z@
	\fi}
% The two earlier attempts, kept for reference. Both ended the description row
% through a hook, \cv@descrowend, which chose between them:
%   (1) \\*[\cvtagsep] -- it only deepens the row's strut, and a description
%       minipage is far deeper than any strut, so under a bullet list it
%       added nothing. It showed only under an empty description.
%   (2) \\*\noalign{\nobreak\vskip\cvtagsep} -- the space appeared, but a
%       \noalign skip sits between rows, so the timeline's vertical line
%       broke at every tag row (seen in the render).
% \newcommand\cv@descrowend{\ifdim\cvtagsep>\z@
% 		\expandafter\cv@descrowend@sep
% 	\else
% 		\expandafter\cv@descrowend@plain
% 	\fi}
% \newcommand\cv@descrowend@sep{\\*\noalign{\nobreak\vskip\cvtagsep}}
% \newcommand\cv@descrowend@plain{\\*}
% The tag row expands its argument once before \foreach sees it (2026-09-23).
% \foreach takes a braced list literally, so {\cvFhgTags} would come out as a
% single tag holding the whole list. A literal list starts with a character,
% which \expandafter leaves alone, so every existing call typesets exactly as
% before -- and no tag list in data/cv starts with a macro. The row before,
% kept:
% 	& \footnotesize{\foreach \n in {#7}{\cvtag{\n}}} \\
\newcommand\cv@tagrow[1]{\foreach \n in {#1}{\cvtag{\n}}}
\renewcommand\experience[7]{
	\textbf{#1}    & \textbf{#2, \textsc{#3}, #4}   			   \\*
	\textbf{#5}    & \begin{minipage}[t]{\rightcolumnlength}
		#6 \cv@descpad
	\end{minipage}								   \\*
	& \footnotesize{\expandafter\cv@tagrow\expandafter{#7}} \\
}
\makeatother

\providecommand{\cvprofile}{default}
\input{profiles/\cvprofile.tex}

% The word after the name in the footer: "CV" here, redefined by a cover
% letter. The CV's own output is unchanged.
\providecommand{\cvfootertitle}{CV}

% Cover letters (2026-09-23). A letter build passes \def\cvletter{<name>} the
% way \cvprofile is passed (scripts/build_in_container.sh, target `letter`),
% and <name> is the application's CV profile, which has already been read
% above. So the letter gets this document's class, fonts, colours, header --
% with that profile's tagline and relocation line -- and footer, and cannot
% drift from its CV. letters/<name>.tex redefines \cvBody as the letter and
% \cvfootertitle. brand-letter.tex, next to the class, holds the letter
% layout. A CV build never defines \cvletter, so nothing here runs for it.
\ifdefined\cvletter
	\input{brand-letter.tex}
	\input{letters/\cvletter.tex}
\fi

\name{\brandFirstName}{\brandLastName}
\tagline{\cvTagline}%|
% No \photo: dropping it makes \makecvheader give the header block the full
% \linewidth, which is what lets the three columns below span the page. The
% Gauss citation, which used to occupy the photo slot, is now column three.
%
% Contact details stacked in one column ran nine lines deep and set the height
% of the whole header. Split across two columns beside the citation, the block
% is five lines instead -- the saving the one-page layout needs.
%
% Each entry is an \mbox (see \socialtext in the class), so it cannot wrap: a
% column narrower than its longest entry overflows rather than reflowing.
% Widths are therefore sized to the longest line in each column, not evenly.
%
% The class appends \hspace{1em} to every entry to separate entries that share
% a line. Here each entry owns its line, so that 1em is dead width at the right
% edge of every column -- and it was exactly what pushed the longest entries
% past their column. Dropped locally rather than in the class, which other
% documents share.
%
% The icon sits in a fixed-width box rather than being followed by a fixed
% space, so every entry's text starts at the same x no matter how wide its
% symbol is. With a plain \hspace the text began wherever the glyph happened
% to end: the FontAwesome symbols here differ by up to 1.9pt (map marker vs.
% info), so the middle column stepped in and out by that much down its five
% lines. It had been patched per line with hand-tuned \hspace values in the
% entries below -- 0.3em here, 0.5em there, 0.4em on the phone -- which is
% both unmaintainable and was simply missing on the citizenship line, leaving
% it 4.5pt left of its neighbours. Those kerns are gone; do not add new ones.
%
% 1.45em is the widest symbol in this header (the link glyph, 8.33pt at this
% size) plus the 0.5em gap the class used. It costs the widest entry 0.64pt of
% extra width -- the columns are sized to their longest line, so check the
% German build for an overfull box after changing it.
\newcommand{\socialiconwidth}{1.45em}
\renewcommand{\socialtext}[2]{\mbox{\makebox[\socialiconwidth][l]{\textcolor{symbolcolor}{#1}}#2}}
\renewcommand{\sociallink}[3]{\mbox{\makebox[\socialiconwidth][l]{\textcolor{symbolcolor}{#1}}\link{#2}{#3}}}

% Handle instead of full URL for these two: their "linkedin.com/in/" and
% "youtube.com/@" prefixes are already said by the icon beside them, and they
% were the two widest unbreakable entries in the block. Dropping the prefixes
% is what leaves the group narrower than \linewidth -- without it there is no
% spare width and "centred" would be indistinguishable from full-width. The
% links themselves still point at the full URLs.
\renewcommand*{\linkedin}[1]{\sociallink{\linkedinSymbol}{https://www.linkedin.com/in/#1/}{#1}}
\renewcommand*{\youtube}[1]{\sociallink{\youtubeSymbol}{https://youtube.com/@#1}{@#1}}

% \makebox centres the three columns as one group: the gutters are fixed, so
% the leftover width falls outside the group as equal left/right margins
% instead of being pushed between the columns by \hfill.
\socialinfo{
	\makebox[\linewidth][c]{%
	% Sized to the e-mail address, now the longest entry in this column.
	\begin{minipage}[t]{0.26\linewidth}
		\raggedright
		\githubProfile\newline
		\linkedin{\brandLinkedin}\newline
		\brandWebsite\newline
		% YouTube commented out -- no videos published yet, and a channel handle
		% that leads to an empty channel costs more than it pays. Note the
		% \newline moved up to the e-mail line: this entry was last in the
		% minipage and carried no \newline of its own, so commenting it out
		% alone would have left the e-mail's \newline dangling and set an empty
		% final line. Restore both together.
		\email{\brandEmail}
		%\newline
		%\youtube{\brandYoutube}
	\end{minipage}%
	\hspace{1.5em}%
	% Widest column: it carries the German "Geburtstag: ..., Schwäbisch Hall",
	% which is longer than its English counterpart. Sized for the German build.
	\begin{minipage}[t]{0.395\linewidth}
		\raggedright
		\address{Stuttgart}\newline
		\ifcvrelocate
		\infos{\IfLanguageName{english}{Open to relocating to Paris}{Umzug nach Paris möglich}}\newline
		\fi
		\infos{\IfLanguageName{english}{German Citizen}{Deutscher Staatsbürger}}\newline
		% ", Germany" dropped: the citizenship line above already says it, and
		% this is the longest entry in the column.
		% Geburtsort commented out per request -- the date stays, the place is
		% gone. Restore by putting ", Schwäbisch Hall" back inside both branches:
		%   {Birthday: \DTMdisplaydate{1995}{12}{25}, Schwäbisch Hall}
		%   {Geburtstag: \DTMdisplaydate{1995}{12}{25}, Schwäbisch Hall}
		% This was the widest entry in the column and the one the column width
		% was sized for (see the note above the minipage); with it gone the
		% column has slack, so nothing needs re-measuring until it comes back.
		% The empty fourth argument is required, not decoration: \DTMdisplaydate
		% takes {year}{month}{day}{dow}. The old line got away with three brace
		% groups because the ", Schwäbisch Hall" that followed handed it a comma
		% as the day-of-week, which the date style then printed as nothing. Take
		% the place away and the macro swallows the closing brace instead --
		% "Argument of \DTMdisplaydate has an extra }", and no PDF at all.
		% Date of birth commented out on request -- the birthplace went earlier,
		% now the date follows. Costs the layout nothing: this column drops from
		% five entries to four, and the contact column on the left still sets the
		% header height at five. Restore by uncommenting; note the empty fourth
		% argument to \DTMdisplaydate is required (see the note above).
		%\infos{\IfLanguageName{english}{Birthday: \DTMdisplaydate{1995}{12}{25}{}}{Geburtstag: \DTMdisplaydate{1995}{12}{25}{}}}\newline
		\infos{\IfLanguageName{english}{Driving License: Class B}{Führerschein: Klasse B}}\newline
		\smartphone{+49 15253799950}
	\end{minipage}%
	\hspace{1.5em}%
	\begin{minipage}[t]{0.24\linewidth}
		\raggedright
		\footnotesize\itshape
		\IfLanguageName{english}{%
			\textquote[Carl Friedrich Gauss]{Truly, it is not knowledge, but learning; not possession, but acquisition; not being there, but arriving -- that grants the greatest joy.}%
		}{%
			\textquote[Carl Friedrich Gau\ss]{Wahrlich es ist nicht das Wissen, sondern das Lernen, nicht das Besitzen, sondern das Erwerben, nicht das Da-Seyn, sondern das Hinkommen, was den grössten Genuss gewährt.}%
		}%
	\end{minipage}%
	}%
}

\begin{document}

\makecvheader

\makecvfooter
% The left footer field was empty. The most credible line for "I work in the
% open" is the one showing this PDF is itself the artifact of a public
% pipeline -- and in the footer it costs the one-page layout nothing.
{\footnotesize\textsc{Typeset with}\ \link{\brandGithubUrl/DocumANTation}{DocumANTation}}
% "CV" became \cvfootertitle (2026-09-23), so a cover letter can say what it
% is. The line before, kept:
% {\textsc{\brandName{} - CV}}
{\textsc{\brandName{} - \cvfootertitle}}
{\thepage}

% The section list moved into profiles/: which sections fit is a per-target
% decision, and it was already being made here by commenting lines in and out.
\cvBody

\end{document}
"
% which is what CV_PROFILE in scripts/build_in_container.sh does. Every profile
% defines \cvTagline and \cvBody; see profiles/README.md.
%
% Profiles select *sections*, they do not fork prose: a per-profile branch
% inside every section file is the same trap data/cv/README.md flags for a
% third language. Role-specific wording lives in its own section_headline_*
% file, which the profile inputs instead of the general one.
%
% The one exception to "profiles do not edit sections": the A-levels entry in
% section_education.tex. A profile aimed at a role that requires a Master's
% spends three lines on a 2006 Gymnasium entry it does not need, and the page
% budget is the whole constraint. A \newif rather than an \ifdefstring: the
% entry it guards contains blank lines, which an argument-grabbing conditional
% would choke on. Declared here, before the profile is read, so a profile can
% flip it.
\newif\ifcvshowschool \cvshowschooltrue
% Same idea, one entry up: drops the thesis lines while keeping the degrees,
% their dates, their grades and their subject tags.
\newif\ifcvshowthesislines \cvshowthesislinestrue
% Suppresses the "Teaching and Mentoring" heading so the tutor entry can follow
% the other jobs under Experiences -- which is what it is, and which is what
% makes an unbroken 2018-today employment line visible on one page. Only useful
% for a profile that inputs section_teaching_mentoring directly after
% section_experience.
\newif\ifcvshowteachingtitle \cvshowteachingtitletrue
% Drops the CUDA-kernels bullet. The one bullet a platform or infrastructure
% posting does not buy, and the cheapest two lines in the Fraunhofer entry for
% a profile that needs room. On for every research-facing profile.
\newif\ifcvshowcudakernels \cvshowcudakernelstrue
% The AI-gateway bullet, and the only flag here that defaults OFF: it is the
% one place the CV claims Go, so a profile opts in rather than every existing
% profile paying two lines it did not budget for. \ifcvrelocate is the same
% shape. `all` turns it on, which is what `all` means.
\newif\ifcvshowgateway \cvshowgatewayfalse
% Drops the RemaNet bullet -- three lines, the longest in the entry. Edge-AI
% product work: central to a research profile, the first thing a platform
% posting trades away.
\newif\ifcvshowremanet \cvshowremanettrue
% The documentation bullet. Default off for the same reason as the gateway:
% new content added for one application, and every other profile is at zero
% slack. `all` turns it on.
\newif\ifcvshowdocs \cvshowdocsfalse
% The quantum-physics GPU-port bullet (Fraunhofer IPA's quantum-computing
% group). Same shape as the gateway/docs pair: new content for one application
% -- the AMD HPC one -- and the published profiles are at zero slack. `all`
% turns it on. The amd-hpc profile trades the RemaNet bullet for this one.
\newif\ifcvshowquantum \cvshowquantumfalse
% The SemProbe one-liner under the ECML PKDD entry in section_private.tex. It
% used to be commented out and is now behind this flag. Default off because it
% costs two \footnotesize lines in Publications, the column that sets the
% height of the bottom row, and the published profiles have zero slack (the
% German mistral-rse build is the one it pushes over). terabase-cv opts in
% because it is the only section line that says "object detectors", and that
% profile's summary sentence on detectors rests on it. `all` turns it on.
\newif\ifcvshowsemprobe \cvshowsemprobefalse
% The Berghof bullet: object detection on an embedded Berghof PLC fitted with
% a Hailo accelerator. Until 2026-09-22 it sat inside the RemaNet bullet, which
% made it read as part of RemaNet. The owner corrected that: it was a separate
% project, with its own full ML lifecycle. Default ON, so every profile that
% showed the PLC still shows it. amd-hpc turns it off: that profile hides
% RemaNet, so it never showed the PLC.
\newif\ifcvshowberghof \cvshowberghoftrue
% The industrial safety project: person detection, 4 cameras at once, their
% inferences batched, full ML lifecycle (owner, 2026-09-22). New content, so
% it is default OFF like gateway/docs/quantum: the published profiles have
% zero slack. `all` and terabase-cv turn it on.
\newif\ifcvshowsafety \cvshowsafetyfalse
% The build-systems bullet: CMake, GStreamer (Meson), LiteRT-LM (Bazel), and
% fixes upstreamed to LLVM and sccache. On everywhere except terabase-cv,
% where its three lines pay for the Berghof and safety bullets. It is the
% bullet a computer-vision posting has the least use for.
\newif\ifcvshowbuilds \cvshowbuildstrue
% The KI-PrOT bullet: visual classification of parts into drag-out classes
% (surface treatment; ki-prot.de, Fraunhofer IPA develops the AI methods).
% The owner built and annotated the dataset and versioned it with DVC
% (owner, 2026-09-22). New content, so default OFF; terabase-cv and `all`
% turn it on.
\newif\ifcvshowkiprot \cvshowkiprotfalse
% ML-lifecycle detail clauses, appended to existing bullets when on: INT8 and
% the customer field-feedback loop on the Berghof PLC, the owner's own
% self-labeling research tool (some of the owner's projects used its generated
% labels -- reworded 2026-09-23; first put as "its labels partly trained the
% models"), and
% MLflow/ClearML experiment tracking. All owner-confirmed on 2026-09-22. Default
% OFF: the base sentences stay exactly as they were, so zero-slack profiles
% are untouched. terabase-cv and `all` turn it on.
\newif\ifcvshowmldetail \cvshowmldetailfalse
% Stack-detail clauses, appended when on (owner, 2026-09-23): the coding
% agent's open-weight models come from Hugging Face, and RemaNet's web part
% used TypeScript and Node.js. Since the same day it also names Docker in the
% CI/CD bullet ("cut Docker builds"): Docker daily at Fraunhofer, advanced
% builds (owner). Default OFF, and when off the base sentences
% typeset exactly as before, so no other profile's lines move. parallel-ai
% and `all` turn it on.
\newif\ifcvshowstackdetail \cvshowstackdetailfalse
% Drops the CAS description, keeping the entry: its dates, title, employer and
% tags all stay, so the timeline is untouched. The ONNX/mixed-reality bullet is
% the one description on the page a platform posting has no use for. Guards the
% whole itemize, not the \item -- an itemize with no \item is a LaTeX error.
\newif\ifcvshowcasdetail \cvshowcasdetailtrue
% Off by default, on for a profile aimed at a job in another city. The header
% is where a reader checks whether the geography works, and it answers the
% question "Stuttgart, but the job is in Paris?" in the same place it is asked.
% Costs nothing: with the birthday and the YouTube handle gone both contact
% columns run four lines, and the citation column beside them is five (English)
% or seven (German), so it is the citation that sets the header height. Only
% grow this past five lines with the German build in front of you.
\newif\ifcvrelocate \cvrelocatefalse
% The order of the Fraunhofer IPA bullets. section_experience.tex defines each
% one as \cvFhg<Name> (each still behind its own \ifcvshow... flag) and prints
% them through this macro, so a profile can lead with what its posting asks
% for first without forking any prose: \renewcommand it in the profile.
% This default is the order the section file has always had. List every
% bullet in a profile's order too, even hidden ones: a bullet left out of the
% list never prints, whatever its flag says, and that failure is silent.
\newcommand{\cvFraunhoferOrder}{%
	\cvFhgKernels\cvFhgQuantum\cvFhgCicd\cvFhgBuilds\cvFhgGateway\cvFhgDocs%
	\cvFhgVlm\cvFhgAgent\cvFhgRemanet\cvFhgBerghof\cvFhgSafety\cvFhgKiprot}
% The Fraunhofer IPA tag row, set per profile the same way (2026-09-23). The
% row is one line wide and every tag in it competes for that line, so a
% profile whose posting has no use for CMake, Bazel and Meson -- the build
% systems of a bullet it hides -- can spend the room on what it does ask for:
% \renewcommand it in the profile. This default is the list the section file
% had before, so no other profile changes. The \experience override below
% expands the tag argument once, which is what lets a macro stand in for a
% literal list there.
\newcommand{\cvFhgTags}{Python, uv, Go, C++, CUDA, PyTorch, CMake, Bazel, Meson, vLLM, LoRA, GitLab CI/CD}
% \cvshowthesislinesfalse hid the thesis text but not the row holding it:
% \education always emits its description row, so switching the flag off
% left an empty minipage -- a blank line under every degree, which is exactly
% the vertical space the flag was set to reclaim. Emitting that row only when
% there is something to put in it makes the flag pay what it promises.
% Overridden here rather than in myCV_METADATA.cls, which other documents
% share, following the \socialtext/\sociallink precedent below.
\renewcommand\education[8]{
	\textbf{#1}    & \textbf{#2, \textsc{#3}, #4}   \\*
	\textbf{#5}    & #6   \\*
	\ifcvshowthesislines
	& \begin{minipage}[t]{\rightcolumnlength}
		#7
	\end{minipage}   \\*
	\fi
	& \footnotesize{\foreach \n in {#8}{\cvtag{\n}}} \\
}
% Extra space between an experience's description and its tag row. The class
% sets the tag boxes directly under the last description line, and when that
% line has descenders the boxes nearly touch it ("much too less space", owner,
% 2026-09-22, on the Fraunhofer entry). \cvtagsep is 0pt by default, and at
% zero the row ends in the class's plain \\* -- token for token the class's own
% definition, so no profile changes unless it sets the length. terabase-cv
% does. \expandafter drops the \else/\fi before \\ runs, so no conditional is
% left open across the row end.
% The class's \experience, copied verbatim below except for the one row end.
% Overridden here, not in myCV_METADATA.cls, for the same reason as
% \education above.
\newlength{\cvtagsep}
\setlength{\cvtagsep}{0pt}
\makeatletter
% How it works (third attempt, 2026-09-22): a zero-width rule of height
% \cvtagsep closes the description minipage. It deepens that cell by exactly
% \cvtagsep, the timeline's vertical line stays unbroken (same row), and an
% empty description (a hidden CAS detail) gains nothing: there the rule is
% the minipage's first item and only sets its height, which the row strut
% already exceeds. At 0pt it expands to nothing, so the class layout is
% unchanged token for token.
\newcommand\cv@descpad{\ifdim\cvtagsep>\z@
		\par\nointerlineskip\hrule\@height\cvtagsep\@width\z@
	\fi}
% The two earlier attempts, kept for reference. Both ended the description row
% through a hook, \cv@descrowend, which chose between them:
%   (1) \\*[\cvtagsep] -- it only deepens the row's strut, and a description
%       minipage is far deeper than any strut, so under a bullet list it
%       added nothing. It showed only under an empty description.
%   (2) \\*\noalign{\nobreak\vskip\cvtagsep} -- the space appeared, but a
%       \noalign skip sits between rows, so the timeline's vertical line
%       broke at every tag row (seen in the render).
% \newcommand\cv@descrowend{\ifdim\cvtagsep>\z@
% 		\expandafter\cv@descrowend@sep
% 	\else
% 		\expandafter\cv@descrowend@plain
% 	\fi}
% \newcommand\cv@descrowend@sep{\\*\noalign{\nobreak\vskip\cvtagsep}}
% \newcommand\cv@descrowend@plain{\\*}
% The tag row expands its argument once before \foreach sees it (2026-09-23).
% \foreach takes a braced list literally, so {\cvFhgTags} would come out as a
% single tag holding the whole list. A literal list starts with a character,
% which \expandafter leaves alone, so every existing call typesets exactly as
% before -- and no tag list in data/cv starts with a macro. The row before,
% kept:
% 	& \footnotesize{\foreach \n in {#7}{\cvtag{\n}}} \\
\newcommand\cv@tagrow[1]{\foreach \n in {#1}{\cvtag{\n}}}
\renewcommand\experience[7]{
	\textbf{#1}    & \textbf{#2, \textsc{#3}, #4}   			   \\*
	\textbf{#5}    & \begin{minipage}[t]{\rightcolumnlength}
		#6 \cv@descpad
	\end{minipage}								   \\*
	& \footnotesize{\expandafter\cv@tagrow\expandafter{#7}} \\
}
\makeatother

\providecommand{\cvprofile}{default}
\input{profiles/\cvprofile.tex}

% The word after the name in the footer: "CV" here, redefined by a cover
% letter. The CV's own output is unchanged.
\providecommand{\cvfootertitle}{CV}

% Cover letters (2026-09-23). A letter build passes \def\cvletter{<name>} the
% way \cvprofile is passed (scripts/build_in_container.sh, target `letter`),
% and <name> is the application's CV profile, which has already been read
% above. So the letter gets this document's class, fonts, colours, header --
% with that profile's tagline and relocation line -- and footer, and cannot
% drift from its CV. letters/<name>.tex redefines \cvBody as the letter and
% \cvfootertitle. brand-letter.tex, next to the class, holds the letter
% layout. A CV build never defines \cvletter, so nothing here runs for it.
\ifdefined\cvletter
	% Cover-letter layout for the CV class (myCV_METADATA), added 2026-09-23.
%
% data/cv/cv.tex reads this file only when a letter is built -- \cvletter is
% defined, see scripts/build_in_container.sh, target `letter` -- and then
% reads data/cv/letters/<name>.tex, which calls the macros below. A letter is
% the CV document with a different \cvBody: same class, same brand colours
% and fonts, same header (name, the application's tagline, contact columns),
% same footer. Layout lives here and content lives in data/cv/letters/, the
% same split as the CV itself.
%
% \providecommand throughout, so a letter file may define its own first.
%
% \cvletterhead{<place, date>}{<recipient, lines split by \\>}{<subject>}
%     Recipient on the left, place and date on the right, then the subject
%     in the heading colour.
% \cvletterclosing{<closing phrase>}
%     The closing phrase, room for a handwritten signature, then the name --
%     \brandName, never typed (tests/test_identity_is_single_source.py).

\providecommand{\cvletterhead}[3]{%
	\par\vspace{0.4cm}%
	\noindent
	\begin{minipage}[t]{0.58\linewidth}\raggedright #2\end{minipage}%
	\hfill
	\begin{minipage}[t]{0.38\linewidth}\raggedleft #1\end{minipage}%
	\par\vspace{0.7cm}%
	{\large\color{headingcolor}\textbf{#3}}%
	\par\vspace{0.35cm}%
}

\providecommand{\cvletterclosing}[1]{%
	\par\vspace{0.2cm}%
	#1,\par
	\vspace{0.7cm}%
	{\color{headingcolor}\textbf{\brandName}}\par
}

	\input{letters/\cvletter.tex}
\fi

\name{\brandFirstName}{\brandLastName}
\tagline{\cvTagline}%|
% No \photo: dropping it makes \makecvheader give the header block the full
% \linewidth, which is what lets the three columns below span the page. The
% Gauss citation, which used to occupy the photo slot, is now column three.
%
% Contact details stacked in one column ran nine lines deep and set the height
% of the whole header. Split across two columns beside the citation, the block
% is five lines instead -- the saving the one-page layout needs.
%
% Each entry is an \mbox (see \socialtext in the class), so it cannot wrap: a
% column narrower than its longest entry overflows rather than reflowing.
% Widths are therefore sized to the longest line in each column, not evenly.
%
% The class appends \hspace{1em} to every entry to separate entries that share
% a line. Here each entry owns its line, so that 1em is dead width at the right
% edge of every column -- and it was exactly what pushed the longest entries
% past their column. Dropped locally rather than in the class, which other
% documents share.
%
% The icon sits in a fixed-width box rather than being followed by a fixed
% space, so every entry's text starts at the same x no matter how wide its
% symbol is. With a plain \hspace the text began wherever the glyph happened
% to end: the FontAwesome symbols here differ by up to 1.9pt (map marker vs.
% info), so the middle column stepped in and out by that much down its five
% lines. It had been patched per line with hand-tuned \hspace values in the
% entries below -- 0.3em here, 0.5em there, 0.4em on the phone -- which is
% both unmaintainable and was simply missing on the citizenship line, leaving
% it 4.5pt left of its neighbours. Those kerns are gone; do not add new ones.
%
% 1.45em is the widest symbol in this header (the link glyph, 8.33pt at this
% size) plus the 0.5em gap the class used. It costs the widest entry 0.64pt of
% extra width -- the columns are sized to their longest line, so check the
% German build for an overfull box after changing it.
\newcommand{\socialiconwidth}{1.45em}
\renewcommand{\socialtext}[2]{\mbox{\makebox[\socialiconwidth][l]{\textcolor{symbolcolor}{#1}}#2}}
\renewcommand{\sociallink}[3]{\mbox{\makebox[\socialiconwidth][l]{\textcolor{symbolcolor}{#1}}\link{#2}{#3}}}

% Handle instead of full URL for these two: their "linkedin.com/in/" and
% "youtube.com/@" prefixes are already said by the icon beside them, and they
% were the two widest unbreakable entries in the block. Dropping the prefixes
% is what leaves the group narrower than \linewidth -- without it there is no
% spare width and "centred" would be indistinguishable from full-width. The
% links themselves still point at the full URLs.
\renewcommand*{\linkedin}[1]{\sociallink{\linkedinSymbol}{https://www.linkedin.com/in/#1/}{#1}}
\renewcommand*{\youtube}[1]{\sociallink{\youtubeSymbol}{https://youtube.com/@#1}{@#1}}

% \makebox centres the three columns as one group: the gutters are fixed, so
% the leftover width falls outside the group as equal left/right margins
% instead of being pushed between the columns by \hfill.
\socialinfo{
	\makebox[\linewidth][c]{%
	% Sized to the e-mail address, now the longest entry in this column.
	\begin{minipage}[t]{0.26\linewidth}
		\raggedright
		\githubProfile\newline
		\linkedin{\brandLinkedin}\newline
		\brandWebsite\newline
		% YouTube commented out -- no videos published yet, and a channel handle
		% that leads to an empty channel costs more than it pays. Note the
		% \newline moved up to the e-mail line: this entry was last in the
		% minipage and carried no \newline of its own, so commenting it out
		% alone would have left the e-mail's \newline dangling and set an empty
		% final line. Restore both together.
		\email{\brandEmail}
		%\newline
		%\youtube{\brandYoutube}
	\end{minipage}%
	\hspace{1.5em}%
	% Widest column: it carries the German "Geburtstag: ..., Schwäbisch Hall",
	% which is longer than its English counterpart. Sized for the German build.
	\begin{minipage}[t]{0.395\linewidth}
		\raggedright
		\address{Stuttgart}\newline
		\ifcvrelocate
		\infos{\IfLanguageName{english}{Open to relocating to Paris}{Umzug nach Paris möglich}}\newline
		\fi
		\infos{\IfLanguageName{english}{German Citizen}{Deutscher Staatsbürger}}\newline
		% ", Germany" dropped: the citizenship line above already says it, and
		% this is the longest entry in the column.
		% Geburtsort commented out per request -- the date stays, the place is
		% gone. Restore by putting ", Schwäbisch Hall" back inside both branches:
		%   {Birthday: \DTMdisplaydate{1995}{12}{25}, Schwäbisch Hall}
		%   {Geburtstag: \DTMdisplaydate{1995}{12}{25}, Schwäbisch Hall}
		% This was the widest entry in the column and the one the column width
		% was sized for (see the note above the minipage); with it gone the
		% column has slack, so nothing needs re-measuring until it comes back.
		% The empty fourth argument is required, not decoration: \DTMdisplaydate
		% takes {year}{month}{day}{dow}. The old line got away with three brace
		% groups because the ", Schwäbisch Hall" that followed handed it a comma
		% as the day-of-week, which the date style then printed as nothing. Take
		% the place away and the macro swallows the closing brace instead --
		% "Argument of \DTMdisplaydate has an extra }", and no PDF at all.
		% Date of birth commented out on request -- the birthplace went earlier,
		% now the date follows. Costs the layout nothing: this column drops from
		% five entries to four, and the contact column on the left still sets the
		% header height at five. Restore by uncommenting; note the empty fourth
		% argument to \DTMdisplaydate is required (see the note above).
		%\infos{\IfLanguageName{english}{Birthday: \DTMdisplaydate{1995}{12}{25}{}}{Geburtstag: \DTMdisplaydate{1995}{12}{25}{}}}\newline
		\infos{\IfLanguageName{english}{Driving License: Class B}{Führerschein: Klasse B}}\newline
		\smartphone{+49 15253799950}
	\end{minipage}%
	\hspace{1.5em}%
	\begin{minipage}[t]{0.24\linewidth}
		\raggedright
		\footnotesize\itshape
		\IfLanguageName{english}{%
			\textquote[Carl Friedrich Gauss]{Truly, it is not knowledge, but learning; not possession, but acquisition; not being there, but arriving -- that grants the greatest joy.}%
		}{%
			\textquote[Carl Friedrich Gau\ss]{Wahrlich es ist nicht das Wissen, sondern das Lernen, nicht das Besitzen, sondern das Erwerben, nicht das Da-Seyn, sondern das Hinkommen, was den grössten Genuss gewährt.}%
		}%
	\end{minipage}%
	}%
}

\begin{document}

\makecvheader

\makecvfooter
% The left footer field was empty. The most credible line for "I work in the
% open" is the one showing this PDF is itself the artifact of a public
% pipeline -- and in the footer it costs the one-page layout nothing.
{\footnotesize\textsc{Typeset with}\ \link{\brandGithubUrl/DocumANTation}{DocumANTation}}
% "CV" became \cvfootertitle (2026-09-23), so a cover letter can say what it
% is. The line before, kept:
% {\textsc{\brandName{} - CV}}
{\textsc{\brandName{} - \cvfootertitle}}
{\thepage}

% The section list moved into profiles/: which sections fit is a per-target
% decision, and it was already being made here by commenting lines in and out.
\cvBody

\end{document}
"
% which is what CV_PROFILE in scripts/build_in_container.sh does. Every profile
% defines \cvTagline and \cvBody; see profiles/README.md.
%
% Profiles select *sections*, they do not fork prose: a per-profile branch
% inside every section file is the same trap data/cv/README.md flags for a
% third language. Role-specific wording lives in its own section_headline_*
% file, which the profile inputs instead of the general one.
%
% The one exception to "profiles do not edit sections": the A-levels entry in
% section_education.tex. A profile aimed at a role that requires a Master's
% spends three lines on a 2006 Gymnasium entry it does not need, and the page
% budget is the whole constraint. A \newif rather than an \ifdefstring: the
% entry it guards contains blank lines, which an argument-grabbing conditional
% would choke on. Declared here, before the profile is read, so a profile can
% flip it.
\newif\ifcvshowschool \cvshowschooltrue
% Same idea, one entry up: drops the thesis lines while keeping the degrees,
% their dates, their grades and their subject tags.
\newif\ifcvshowthesislines \cvshowthesislinestrue
% Suppresses the "Teaching and Mentoring" heading so the tutor entry can follow
% the other jobs under Experiences -- which is what it is, and which is what
% makes an unbroken 2018-today employment line visible on one page. Only useful
% for a profile that inputs section_teaching_mentoring directly after
% section_experience.
\newif\ifcvshowteachingtitle \cvshowteachingtitletrue
% Drops the CUDA-kernels bullet. The one bullet a platform or infrastructure
% posting does not buy, and the cheapest two lines in the Fraunhofer entry for
% a profile that needs room. On for every research-facing profile.
\newif\ifcvshowcudakernels \cvshowcudakernelstrue
% The AI-gateway bullet, and the only flag here that defaults OFF: it is the
% one place the CV claims Go, so a profile opts in rather than every existing
% profile paying two lines it did not budget for. \ifcvrelocate is the same
% shape. `all` turns it on, which is what `all` means.
\newif\ifcvshowgateway \cvshowgatewayfalse
% Drops the RemaNet bullet -- three lines, the longest in the entry. Edge-AI
% product work: central to a research profile, the first thing a platform
% posting trades away.
\newif\ifcvshowremanet \cvshowremanettrue
% The documentation bullet. Default off for the same reason as the gateway:
% new content added for one application, and every other profile is at zero
% slack. `all` turns it on.
\newif\ifcvshowdocs \cvshowdocsfalse
% The quantum-physics GPU-port bullet (Fraunhofer IPA's quantum-computing
% group). Same shape as the gateway/docs pair: new content for one application
% -- the AMD HPC one -- and the published profiles are at zero slack. `all`
% turns it on. The amd-hpc profile trades the RemaNet bullet for this one.
\newif\ifcvshowquantum \cvshowquantumfalse
% The SemProbe one-liner under the ECML PKDD entry in section_private.tex. It
% used to be commented out and is now behind this flag. Default off because it
% costs two \footnotesize lines in Publications, the column that sets the
% height of the bottom row, and the published profiles have zero slack (the
% German mistral-rse build is the one it pushes over). terabase-cv opts in
% because it is the only section line that says "object detectors", and that
% profile's summary sentence on detectors rests on it. `all` turns it on.
\newif\ifcvshowsemprobe \cvshowsemprobefalse
% The Berghof bullet: object detection on an embedded Berghof PLC fitted with
% a Hailo accelerator. Until 2026-09-22 it sat inside the RemaNet bullet, which
% made it read as part of RemaNet. The owner corrected that: it was a separate
% project, with its own full ML lifecycle. Default ON, so every profile that
% showed the PLC still shows it. amd-hpc turns it off: that profile hides
% RemaNet, so it never showed the PLC.
\newif\ifcvshowberghof \cvshowberghoftrue
% The industrial safety project: person detection, 4 cameras at once, their
% inferences batched, full ML lifecycle (owner, 2026-09-22). New content, so
% it is default OFF like gateway/docs/quantum: the published profiles have
% zero slack. `all` and terabase-cv turn it on.
\newif\ifcvshowsafety \cvshowsafetyfalse
% The build-systems bullet: CMake, GStreamer (Meson), LiteRT-LM (Bazel), and
% fixes upstreamed to LLVM and sccache. On everywhere except terabase-cv,
% where its three lines pay for the Berghof and safety bullets. It is the
% bullet a computer-vision posting has the least use for.
\newif\ifcvshowbuilds \cvshowbuildstrue
% The KI-PrOT bullet: visual classification of parts into drag-out classes
% (surface treatment; ki-prot.de, Fraunhofer IPA develops the AI methods).
% The owner built and annotated the dataset and versioned it with DVC
% (owner, 2026-09-22). New content, so default OFF; terabase-cv and `all`
% turn it on.
\newif\ifcvshowkiprot \cvshowkiprotfalse
% ML-lifecycle detail clauses, appended to existing bullets when on: INT8 and
% the customer field-feedback loop on the Berghof PLC, the owner's own
% self-labeling research tool (some of the owner's projects used its generated
% labels -- reworded 2026-09-23; first put as "its labels partly trained the
% models"), and
% MLflow/ClearML experiment tracking. All owner-confirmed on 2026-09-22. Default
% OFF: the base sentences stay exactly as they were, so zero-slack profiles
% are untouched. terabase-cv and `all` turn it on.
\newif\ifcvshowmldetail \cvshowmldetailfalse
% Stack-detail clauses, appended when on (owner, 2026-09-23): the coding
% agent's open-weight models come from Hugging Face, and RemaNet's web part
% used TypeScript and Node.js. Since the same day it also names Docker in the
% CI/CD bullet ("cut Docker builds"): Docker daily at Fraunhofer, advanced
% builds (owner). Default OFF, and when off the base sentences
% typeset exactly as before, so no other profile's lines move. parallel-ai
% and `all` turn it on.
\newif\ifcvshowstackdetail \cvshowstackdetailfalse
% Drops the CAS description, keeping the entry: its dates, title, employer and
% tags all stay, so the timeline is untouched. The ONNX/mixed-reality bullet is
% the one description on the page a platform posting has no use for. Guards the
% whole itemize, not the \item -- an itemize with no \item is a LaTeX error.
\newif\ifcvshowcasdetail \cvshowcasdetailtrue
% Off by default, on for a profile aimed at a job in another city. The header
% is where a reader checks whether the geography works, and it answers the
% question "Stuttgart, but the job is in Paris?" in the same place it is asked.
% Costs nothing: with the birthday and the YouTube handle gone both contact
% columns run four lines, and the citation column beside them is five (English)
% or seven (German), so it is the citation that sets the header height. Only
% grow this past five lines with the German build in front of you.
\newif\ifcvrelocate \cvrelocatefalse
% The order of the Fraunhofer IPA bullets. section_experience.tex defines each
% one as \cvFhg<Name> (each still behind its own \ifcvshow... flag) and prints
% them through this macro, so a profile can lead with what its posting asks
% for first without forking any prose: \renewcommand it in the profile.
% This default is the order the section file has always had. List every
% bullet in a profile's order too, even hidden ones: a bullet left out of the
% list never prints, whatever its flag says, and that failure is silent.
\newcommand{\cvFraunhoferOrder}{%
	\cvFhgKernels\cvFhgQuantum\cvFhgCicd\cvFhgBuilds\cvFhgGateway\cvFhgDocs%
	\cvFhgVlm\cvFhgAgent\cvFhgRemanet\cvFhgBerghof\cvFhgSafety\cvFhgKiprot}
% The Fraunhofer IPA tag row, set per profile the same way (2026-09-23). The
% row is one line wide and every tag in it competes for that line, so a
% profile whose posting has no use for CMake, Bazel and Meson -- the build
% systems of a bullet it hides -- can spend the room on what it does ask for:
% \renewcommand it in the profile. This default is the list the section file
% had before, so no other profile changes. The \experience override below
% expands the tag argument once, which is what lets a macro stand in for a
% literal list there.
\newcommand{\cvFhgTags}{Python, uv, Go, C++, CUDA, PyTorch, CMake, Bazel, Meson, vLLM, LoRA, GitLab CI/CD}
% \cvshowthesislinesfalse hid the thesis text but not the row holding it:
% \education always emits its description row, so switching the flag off
% left an empty minipage -- a blank line under every degree, which is exactly
% the vertical space the flag was set to reclaim. Emitting that row only when
% there is something to put in it makes the flag pay what it promises.
% Overridden here rather than in myCV_METADATA.cls, which other documents
% share, following the \socialtext/\sociallink precedent below.
\renewcommand\education[8]{
	\textbf{#1}    & \textbf{#2, \textsc{#3}, #4}   \\*
	\textbf{#5}    & #6   \\*
	\ifcvshowthesislines
	& \begin{minipage}[t]{\rightcolumnlength}
		#7
	\end{minipage}   \\*
	\fi
	& \footnotesize{\foreach \n in {#8}{\cvtag{\n}}} \\
}
% Extra space between an experience's description and its tag row. The class
% sets the tag boxes directly under the last description line, and when that
% line has descenders the boxes nearly touch it ("much too less space", owner,
% 2026-09-22, on the Fraunhofer entry). \cvtagsep is 0pt by default, and at
% zero the row ends in the class's plain \\* -- token for token the class's own
% definition, so no profile changes unless it sets the length. terabase-cv
% does. \expandafter drops the \else/\fi before \\ runs, so no conditional is
% left open across the row end.
% The class's \experience, copied verbatim below except for the one row end.
% Overridden here, not in myCV_METADATA.cls, for the same reason as
% \education above.
\newlength{\cvtagsep}
\setlength{\cvtagsep}{0pt}
\makeatletter
% How it works (third attempt, 2026-09-22): a zero-width rule of height
% \cvtagsep closes the description minipage. It deepens that cell by exactly
% \cvtagsep, the timeline's vertical line stays unbroken (same row), and an
% empty description (a hidden CAS detail) gains nothing: there the rule is
% the minipage's first item and only sets its height, which the row strut
% already exceeds. At 0pt it expands to nothing, so the class layout is
% unchanged token for token.
\newcommand\cv@descpad{\ifdim\cvtagsep>\z@
		\par\nointerlineskip\hrule\@height\cvtagsep\@width\z@
	\fi}
% The two earlier attempts, kept for reference. Both ended the description row
% through a hook, \cv@descrowend, which chose between them:
%   (1) \\*[\cvtagsep] -- it only deepens the row's strut, and a description
%       minipage is far deeper than any strut, so under a bullet list it
%       added nothing. It showed only under an empty description.
%   (2) \\*\noalign{\nobreak\vskip\cvtagsep} -- the space appeared, but a
%       \noalign skip sits between rows, so the timeline's vertical line
%       broke at every tag row (seen in the render).
% \newcommand\cv@descrowend{\ifdim\cvtagsep>\z@
% 		\expandafter\cv@descrowend@sep
% 	\else
% 		\expandafter\cv@descrowend@plain
% 	\fi}
% \newcommand\cv@descrowend@sep{\\*\noalign{\nobreak\vskip\cvtagsep}}
% \newcommand\cv@descrowend@plain{\\*}
% The tag row expands its argument once before \foreach sees it (2026-09-23).
% \foreach takes a braced list literally, so {\cvFhgTags} would come out as a
% single tag holding the whole list. A literal list starts with a character,
% which \expandafter leaves alone, so every existing call typesets exactly as
% before -- and no tag list in data/cv starts with a macro. The row before,
% kept:
% 	& \footnotesize{\foreach \n in {#7}{\cvtag{\n}}} \\
\newcommand\cv@tagrow[1]{\foreach \n in {#1}{\cvtag{\n}}}
\renewcommand\experience[7]{
	\textbf{#1}    & \textbf{#2, \textsc{#3}, #4}   			   \\*
	\textbf{#5}    & \begin{minipage}[t]{\rightcolumnlength}
		#6 \cv@descpad
	\end{minipage}								   \\*
	& \footnotesize{\expandafter\cv@tagrow\expandafter{#7}} \\
}
\makeatother

\providecommand{\cvprofile}{default}
\input{profiles/\cvprofile.tex}

% The word after the name in the footer: "CV" here, redefined by a cover
% letter. The CV's own output is unchanged.
\providecommand{\cvfootertitle}{CV}

% Cover letters (2026-09-23). A letter build passes \def\cvletter{<name>} the
% way \cvprofile is passed (scripts/build_in_container.sh, target `letter`),
% and <name> is the application's CV profile, which has already been read
% above. So the letter gets this document's class, fonts, colours, header --
% with that profile's tagline and relocation line -- and footer, and cannot
% drift from its CV. letters/<name>.tex redefines \cvBody as the letter and
% \cvfootertitle. brand-letter.tex, next to the class, holds the letter
% layout. A CV build never defines \cvletter, so nothing here runs for it.
\ifdefined\cvletter
	% Cover-letter layout for the CV class (myCV_METADATA), added 2026-09-23.
%
% data/cv/cv.tex reads this file only when a letter is built -- \cvletter is
% defined, see scripts/build_in_container.sh, target `letter` -- and then
% reads data/cv/letters/<name>.tex, which calls the macros below. A letter is
% the CV document with a different \cvBody: same class, same brand colours
% and fonts, same header (name, the application's tagline, contact columns),
% same footer. Layout lives here and content lives in data/cv/letters/, the
% same split as the CV itself.
%
% \providecommand throughout, so a letter file may define its own first.
%
% \cvletterhead{<place, date>}{<recipient, lines split by \\>}{<subject>}
%     Recipient on the left, place and date on the right, then the subject
%     in the heading colour.
% \cvletterclosing{<closing phrase>}
%     The closing phrase, room for a handwritten signature, then the name --
%     \brandName, never typed (tests/test_identity_is_single_source.py).

\providecommand{\cvletterhead}[3]{%
	\par\vspace{0.4cm}%
	\noindent
	\begin{minipage}[t]{0.58\linewidth}\raggedright #2\end{minipage}%
	\hfill
	\begin{minipage}[t]{0.38\linewidth}\raggedleft #1\end{minipage}%
	\par\vspace{0.7cm}%
	{\large\color{headingcolor}\textbf{#3}}%
	\par\vspace{0.35cm}%
}

\providecommand{\cvletterclosing}[1]{%
	\par\vspace{0.2cm}%
	#1,\par
	\vspace{0.7cm}%
	{\color{headingcolor}\textbf{\brandName}}\par
}

	\input{letters/\cvletter.tex}
\fi

\name{\brandFirstName}{\brandLastName}
\tagline{\cvTagline}%|
% No \photo: dropping it makes \makecvheader give the header block the full
% \linewidth, which is what lets the three columns below span the page. The
% Gauss citation, which used to occupy the photo slot, is now column three.
%
% Contact details stacked in one column ran nine lines deep and set the height
% of the whole header. Split across two columns beside the citation, the block
% is five lines instead -- the saving the one-page layout needs.
%
% Each entry is an \mbox (see \socialtext in the class), so it cannot wrap: a
% column narrower than its longest entry overflows rather than reflowing.
% Widths are therefore sized to the longest line in each column, not evenly.
%
% The class appends \hspace{1em} to every entry to separate entries that share
% a line. Here each entry owns its line, so that 1em is dead width at the right
% edge of every column -- and it was exactly what pushed the longest entries
% past their column. Dropped locally rather than in the class, which other
% documents share.
%
% The icon sits in a fixed-width box rather than being followed by a fixed
% space, so every entry's text starts at the same x no matter how wide its
% symbol is. With a plain \hspace the text began wherever the glyph happened
% to end: the FontAwesome symbols here differ by up to 1.9pt (map marker vs.
% info), so the middle column stepped in and out by that much down its five
% lines. It had been patched per line with hand-tuned \hspace values in the
% entries below -- 0.3em here, 0.5em there, 0.4em on the phone -- which is
% both unmaintainable and was simply missing on the citizenship line, leaving
% it 4.5pt left of its neighbours. Those kerns are gone; do not add new ones.
%
% 1.45em is the widest symbol in this header (the link glyph, 8.33pt at this
% size) plus the 0.5em gap the class used. It costs the widest entry 0.64pt of
% extra width -- the columns are sized to their longest line, so check the
% German build for an overfull box after changing it.
\newcommand{\socialiconwidth}{1.45em}
\renewcommand{\socialtext}[2]{\mbox{\makebox[\socialiconwidth][l]{\textcolor{symbolcolor}{#1}}#2}}
\renewcommand{\sociallink}[3]{\mbox{\makebox[\socialiconwidth][l]{\textcolor{symbolcolor}{#1}}\link{#2}{#3}}}

% Handle instead of full URL for these two: their "linkedin.com/in/" and
% "youtube.com/@" prefixes are already said by the icon beside them, and they
% were the two widest unbreakable entries in the block. Dropping the prefixes
% is what leaves the group narrower than \linewidth -- without it there is no
% spare width and "centred" would be indistinguishable from full-width. The
% links themselves still point at the full URLs.
\renewcommand*{\linkedin}[1]{\sociallink{\linkedinSymbol}{https://www.linkedin.com/in/#1/}{#1}}
\renewcommand*{\youtube}[1]{\sociallink{\youtubeSymbol}{https://youtube.com/@#1}{@#1}}

% \makebox centres the three columns as one group: the gutters are fixed, so
% the leftover width falls outside the group as equal left/right margins
% instead of being pushed between the columns by \hfill.
\socialinfo{
	\makebox[\linewidth][c]{%
	% Sized to the e-mail address, now the longest entry in this column.
	\begin{minipage}[t]{0.26\linewidth}
		\raggedright
		\githubProfile\newline
		\linkedin{\brandLinkedin}\newline
		\brandWebsite\newline
		% YouTube commented out -- no videos published yet, and a channel handle
		% that leads to an empty channel costs more than it pays. Note the
		% \newline moved up to the e-mail line: this entry was last in the
		% minipage and carried no \newline of its own, so commenting it out
		% alone would have left the e-mail's \newline dangling and set an empty
		% final line. Restore both together.
		\email{\brandEmail}
		%\newline
		%\youtube{\brandYoutube}
	\end{minipage}%
	\hspace{1.5em}%
	% Widest column: it carries the German "Geburtstag: ..., Schwäbisch Hall",
	% which is longer than its English counterpart. Sized for the German build.
	\begin{minipage}[t]{0.395\linewidth}
		\raggedright
		\address{Stuttgart}\newline
		\ifcvrelocate
		\infos{\IfLanguageName{english}{Open to relocating to Paris}{Umzug nach Paris möglich}}\newline
		\fi
		\infos{\IfLanguageName{english}{German Citizen}{Deutscher Staatsbürger}}\newline
		% ", Germany" dropped: the citizenship line above already says it, and
		% this is the longest entry in the column.
		% Geburtsort commented out per request -- the date stays, the place is
		% gone. Restore by putting ", Schwäbisch Hall" back inside both branches:
		%   {Birthday: \DTMdisplaydate{1995}{12}{25}, Schwäbisch Hall}
		%   {Geburtstag: \DTMdisplaydate{1995}{12}{25}, Schwäbisch Hall}
		% This was the widest entry in the column and the one the column width
		% was sized for (see the note above the minipage); with it gone the
		% column has slack, so nothing needs re-measuring until it comes back.
		% The empty fourth argument is required, not decoration: \DTMdisplaydate
		% takes {year}{month}{day}{dow}. The old line got away with three brace
		% groups because the ", Schwäbisch Hall" that followed handed it a comma
		% as the day-of-week, which the date style then printed as nothing. Take
		% the place away and the macro swallows the closing brace instead --
		% "Argument of \DTMdisplaydate has an extra }", and no PDF at all.
		% Date of birth commented out on request -- the birthplace went earlier,
		% now the date follows. Costs the layout nothing: this column drops from
		% five entries to four, and the contact column on the left still sets the
		% header height at five. Restore by uncommenting; note the empty fourth
		% argument to \DTMdisplaydate is required (see the note above).
		%\infos{\IfLanguageName{english}{Birthday: \DTMdisplaydate{1995}{12}{25}{}}{Geburtstag: \DTMdisplaydate{1995}{12}{25}{}}}\newline
		\infos{\IfLanguageName{english}{Driving License: Class B}{Führerschein: Klasse B}}\newline
		\smartphone{+49 15253799950}
	\end{minipage}%
	\hspace{1.5em}%
	\begin{minipage}[t]{0.24\linewidth}
		\raggedright
		\footnotesize\itshape
		\IfLanguageName{english}{%
			\textquote[Carl Friedrich Gauss]{Truly, it is not knowledge, but learning; not possession, but acquisition; not being there, but arriving -- that grants the greatest joy.}%
		}{%
			\textquote[Carl Friedrich Gau\ss]{Wahrlich es ist nicht das Wissen, sondern das Lernen, nicht das Besitzen, sondern das Erwerben, nicht das Da-Seyn, sondern das Hinkommen, was den grössten Genuss gewährt.}%
		}%
	\end{minipage}%
	}%
}

\begin{document}

\makecvheader

\makecvfooter
% The left footer field was empty. The most credible line for "I work in the
% open" is the one showing this PDF is itself the artifact of a public
% pipeline -- and in the footer it costs the one-page layout nothing.
{\footnotesize\textsc{Typeset with}\ \link{\brandGithubUrl/DocumANTation}{DocumANTation}}
% "CV" became \cvfootertitle (2026-09-23), so a cover letter can say what it
% is. The line before, kept:
% {\textsc{\brandName{} - CV}}
{\textsc{\brandName{} - \cvfootertitle}}
{\thepage}

% The section list moved into profiles/: which sections fit is a per-target
% decision, and it was already being made here by commenting lines in and out.
\cvBody

\end{document}
"
% which is what CV_PROFILE in scripts/build_in_container.sh does. Every profile
% defines \cvTagline and \cvBody; see profiles/README.md.
%
% Profiles select *sections*, they do not fork prose: a per-profile branch
% inside every section file is the same trap data/cv/README.md flags for a
% third language. Role-specific wording lives in its own section_headline_*
% file, which the profile inputs instead of the general one.
%
% The one exception to "profiles do not edit sections": the A-levels entry in
% section_education.tex. A profile aimed at a role that requires a Master's
% spends three lines on a 2006 Gymnasium entry it does not need, and the page
% budget is the whole constraint. A \newif rather than an \ifdefstring: the
% entry it guards contains blank lines, which an argument-grabbing conditional
% would choke on. Declared here, before the profile is read, so a profile can
% flip it.
\newif\ifcvshowschool \cvshowschooltrue
% Same idea, one entry up: drops the thesis lines while keeping the degrees,
% their dates, their grades and their subject tags.
\newif\ifcvshowthesislines \cvshowthesislinestrue
% Suppresses the "Teaching and Mentoring" heading so the tutor entry can follow
% the other jobs under Experiences -- which is what it is, and which is what
% makes an unbroken 2018-today employment line visible on one page. Only useful
% for a profile that inputs section_teaching_mentoring directly after
% section_experience.
\newif\ifcvshowteachingtitle \cvshowteachingtitletrue
% Drops the CUDA-kernels bullet. The one bullet a platform or infrastructure
% posting does not buy, and the cheapest two lines in the Fraunhofer entry for
% a profile that needs room. On for every research-facing profile.
\newif\ifcvshowcudakernels \cvshowcudakernelstrue
% The AI-gateway bullet, and the only flag here that defaults OFF: it is the
% one place the CV claims Go, so a profile opts in rather than every existing
% profile paying two lines it did not budget for. \ifcvrelocate is the same
% shape. `all` turns it on, which is what `all` means.
\newif\ifcvshowgateway \cvshowgatewayfalse
% Drops the RemaNet bullet -- three lines, the longest in the entry. Edge-AI
% product work: central to a research profile, the first thing a platform
% posting trades away.
\newif\ifcvshowremanet \cvshowremanettrue
% The documentation bullet. Default off for the same reason as the gateway:
% new content added for one application, and every other profile is at zero
% slack. `all` turns it on.
\newif\ifcvshowdocs \cvshowdocsfalse
% The quantum-physics GPU-port bullet (Fraunhofer IPA's quantum-computing
% group). Same shape as the gateway/docs pair: new content for one application
% -- the AMD HPC one -- and the published profiles are at zero slack. `all`
% turns it on. The amd-hpc profile trades the RemaNet bullet for this one.
\newif\ifcvshowquantum \cvshowquantumfalse
% The SemProbe one-liner under the ECML PKDD entry in section_private.tex. It
% used to be commented out and is now behind this flag. Default off because it
% costs two \footnotesize lines in Publications, the column that sets the
% height of the bottom row, and the published profiles have zero slack (the
% German mistral-rse build is the one it pushes over). terabase-cv opts in
% because it is the only section line that says "object detectors", and that
% profile's summary sentence on detectors rests on it. `all` turns it on.
\newif\ifcvshowsemprobe \cvshowsemprobefalse
% The Berghof bullet: object detection on an embedded Berghof PLC fitted with
% a Hailo accelerator. Until 2026-09-22 it sat inside the RemaNet bullet, which
% made it read as part of RemaNet. The owner corrected that: it was a separate
% project, with its own full ML lifecycle. Default ON, so every profile that
% showed the PLC still shows it. amd-hpc turns it off: that profile hides
% RemaNet, so it never showed the PLC.
\newif\ifcvshowberghof \cvshowberghoftrue
% The industrial safety project: person detection, 4 cameras at once, their
% inferences batched, full ML lifecycle (owner, 2026-09-22). New content, so
% it is default OFF like gateway/docs/quantum: the published profiles have
% zero slack. `all` and terabase-cv turn it on.
\newif\ifcvshowsafety \cvshowsafetyfalse
% The build-systems bullet: CMake, GStreamer (Meson), LiteRT-LM (Bazel), and
% fixes upstreamed to LLVM and sccache. On everywhere except terabase-cv,
% where its three lines pay for the Berghof and safety bullets. It is the
% bullet a computer-vision posting has the least use for.
\newif\ifcvshowbuilds \cvshowbuildstrue
% The KI-PrOT bullet: visual classification of parts into drag-out classes
% (surface treatment; ki-prot.de, Fraunhofer IPA develops the AI methods).
% The owner built and annotated the dataset and versioned it with DVC
% (owner, 2026-09-22). New content, so default OFF; terabase-cv and `all`
% turn it on.
\newif\ifcvshowkiprot \cvshowkiprotfalse
% ML-lifecycle detail clauses, appended to existing bullets when on: INT8 and
% the customer field-feedback loop on the Berghof PLC, the owner's own
% self-labeling research tool (some of the owner's projects used its generated
% labels -- reworded 2026-09-23; first put as "its labels partly trained the
% models"), and
% MLflow/ClearML experiment tracking. All owner-confirmed on 2026-09-22. Default
% OFF: the base sentences stay exactly as they were, so zero-slack profiles
% are untouched. terabase-cv and `all` turn it on.
\newif\ifcvshowmldetail \cvshowmldetailfalse
% Stack-detail clauses, appended when on (owner, 2026-09-23): the coding
% agent's open-weight models come from Hugging Face, and RemaNet's web part
% used TypeScript and Node.js. Since the same day it also names Docker in the
% CI/CD bullet ("cut Docker builds"): Docker daily at Fraunhofer, advanced
% builds (owner). Default OFF, and when off the base sentences
% typeset exactly as before, so no other profile's lines move. parallel-ai
% and `all` turn it on.
\newif\ifcvshowstackdetail \cvshowstackdetailfalse
% Drops the CAS description, keeping the entry: its dates, title, employer and
% tags all stay, so the timeline is untouched. The ONNX/mixed-reality bullet is
% the one description on the page a platform posting has no use for. Guards the
% whole itemize, not the \item -- an itemize with no \item is a LaTeX error.
\newif\ifcvshowcasdetail \cvshowcasdetailtrue
% Off by default, on for a profile aimed at a job in another city. The header
% is where a reader checks whether the geography works, and it answers the
% question "Stuttgart, but the job is in Paris?" in the same place it is asked.
% Costs nothing: with the birthday and the YouTube handle gone both contact
% columns run four lines, and the citation column beside them is five (English)
% or seven (German), so it is the citation that sets the header height. Only
% grow this past five lines with the German build in front of you.
\newif\ifcvrelocate \cvrelocatefalse
% The order of the Fraunhofer IPA bullets. section_experience.tex defines each
% one as \cvFhg<Name> (each still behind its own \ifcvshow... flag) and prints
% them through this macro, so a profile can lead with what its posting asks
% for first without forking any prose: \renewcommand it in the profile.
% This default is the order the section file has always had. List every
% bullet in a profile's order too, even hidden ones: a bullet left out of the
% list never prints, whatever its flag says, and that failure is silent.
\newcommand{\cvFraunhoferOrder}{%
	\cvFhgKernels\cvFhgQuantum\cvFhgCicd\cvFhgBuilds\cvFhgGateway\cvFhgDocs%
	\cvFhgVlm\cvFhgAgent\cvFhgRemanet\cvFhgBerghof\cvFhgSafety\cvFhgKiprot}
% The Fraunhofer IPA tag row, set per profile the same way (2026-09-23). The
% row is one line wide and every tag in it competes for that line, so a
% profile whose posting has no use for CMake, Bazel and Meson -- the build
% systems of a bullet it hides -- can spend the room on what it does ask for:
% \renewcommand it in the profile. This default is the list the section file
% had before, so no other profile changes. The \experience override below
% expands the tag argument once, which is what lets a macro stand in for a
% literal list there.
\newcommand{\cvFhgTags}{Python, uv, Go, C++, CUDA, PyTorch, CMake, Bazel, Meson, vLLM, LoRA, GitLab CI/CD}
% \cvshowthesislinesfalse hid the thesis text but not the row holding it:
% \education always emits its description row, so switching the flag off
% left an empty minipage -- a blank line under every degree, which is exactly
% the vertical space the flag was set to reclaim. Emitting that row only when
% there is something to put in it makes the flag pay what it promises.
% Overridden here rather than in myCV_METADATA.cls, which other documents
% share, following the \socialtext/\sociallink precedent below.
\renewcommand\education[8]{
	\textbf{#1}    & \textbf{#2, \textsc{#3}, #4}   \\*
	\textbf{#5}    & #6   \\*
	\ifcvshowthesislines
	& \begin{minipage}[t]{\rightcolumnlength}
		#7
	\end{minipage}   \\*
	\fi
	& \footnotesize{\foreach \n in {#8}{\cvtag{\n}}} \\
}
% Extra space between an experience's description and its tag row. The class
% sets the tag boxes directly under the last description line, and when that
% line has descenders the boxes nearly touch it ("much too less space", owner,
% 2026-09-22, on the Fraunhofer entry). \cvtagsep is 0pt by default, and at
% zero the row ends in the class's plain \\* -- token for token the class's own
% definition, so no profile changes unless it sets the length. terabase-cv
% does. \expandafter drops the \else/\fi before \\ runs, so no conditional is
% left open across the row end.
% The class's \experience, copied verbatim below except for the one row end.
% Overridden here, not in myCV_METADATA.cls, for the same reason as
% \education above.
\newlength{\cvtagsep}
\setlength{\cvtagsep}{0pt}
\makeatletter
% How it works (third attempt, 2026-09-22): a zero-width rule of height
% \cvtagsep closes the description minipage. It deepens that cell by exactly
% \cvtagsep, the timeline's vertical line stays unbroken (same row), and an
% empty description (a hidden CAS detail) gains nothing: there the rule is
% the minipage's first item and only sets its height, which the row strut
% already exceeds. At 0pt it expands to nothing, so the class layout is
% unchanged token for token.
\newcommand\cv@descpad{\ifdim\cvtagsep>\z@
		\par\nointerlineskip\hrule\@height\cvtagsep\@width\z@
	\fi}
% The two earlier attempts, kept for reference. Both ended the description row
% through a hook, \cv@descrowend, which chose between them:
%   (1) \\*[\cvtagsep] -- it only deepens the row's strut, and a description
%       minipage is far deeper than any strut, so under a bullet list it
%       added nothing. It showed only under an empty description.
%   (2) \\*\noalign{\nobreak\vskip\cvtagsep} -- the space appeared, but a
%       \noalign skip sits between rows, so the timeline's vertical line
%       broke at every tag row (seen in the render).
% \newcommand\cv@descrowend{\ifdim\cvtagsep>\z@
% 		\expandafter\cv@descrowend@sep
% 	\else
% 		\expandafter\cv@descrowend@plain
% 	\fi}
% \newcommand\cv@descrowend@sep{\\*\noalign{\nobreak\vskip\cvtagsep}}
% \newcommand\cv@descrowend@plain{\\*}
% The tag row expands its argument once before \foreach sees it (2026-09-23).
% \foreach takes a braced list literally, so {\cvFhgTags} would come out as a
% single tag holding the whole list. A literal list starts with a character,
% which \expandafter leaves alone, so every existing call typesets exactly as
% before -- and no tag list in data/cv starts with a macro. The row before,
% kept:
% 	& \footnotesize{\foreach \n in {#7}{\cvtag{\n}}} \\
\newcommand\cv@tagrow[1]{\foreach \n in {#1}{\cvtag{\n}}}
\renewcommand\experience[7]{
	\textbf{#1}    & \textbf{#2, \textsc{#3}, #4}   			   \\*
	\textbf{#5}    & \begin{minipage}[t]{\rightcolumnlength}
		#6 \cv@descpad
	\end{minipage}								   \\*
	& \footnotesize{\expandafter\cv@tagrow\expandafter{#7}} \\
}
\makeatother

\providecommand{\cvprofile}{default}
\input{profiles/\cvprofile.tex}

% The word after the name in the footer: "CV" here, redefined by a cover
% letter. The CV's own output is unchanged.
\providecommand{\cvfootertitle}{CV}

% Cover letters (2026-09-23). A letter build passes \def\cvletter{<name>} the
% way \cvprofile is passed (scripts/build_in_container.sh, target `letter`),
% and <name> is the application's CV profile, which has already been read
% above. So the letter gets this document's class, fonts, colours, header --
% with that profile's tagline and relocation line -- and footer, and cannot
% drift from its CV. letters/<name>.tex redefines \cvBody as the letter and
% \cvfootertitle. brand-letter.tex, next to the class, holds the letter
% layout. A CV build never defines \cvletter, so nothing here runs for it.
\ifdefined\cvletter
	% Cover-letter layout for the CV class (myCV_METADATA), added 2026-09-23.
%
% data/cv/cv.tex reads this file only when a letter is built -- \cvletter is
% defined, see scripts/build_in_container.sh, target `letter` -- and then
% reads data/cv/letters/<name>.tex, which calls the macros below. A letter is
% the CV document with a different \cvBody: same class, same brand colours
% and fonts, same header (name, the application's tagline, contact columns),
% same footer. Layout lives here and content lives in data/cv/letters/, the
% same split as the CV itself.
%
% \providecommand throughout, so a letter file may define its own first.
%
% \cvletterhead{<place, date>}{<recipient, lines split by \\>}{<subject>}
%     Recipient on the left, place and date on the right, then the subject
%     in the heading colour.
% \cvletterclosing{<closing phrase>}
%     The closing phrase, room for a handwritten signature, then the name --
%     \brandName, never typed (tests/test_identity_is_single_source.py).

\providecommand{\cvletterhead}[3]{%
	\par\vspace{0.4cm}%
	\noindent
	\begin{minipage}[t]{0.58\linewidth}\raggedright #2\end{minipage}%
	\hfill
	\begin{minipage}[t]{0.38\linewidth}\raggedleft #1\end{minipage}%
	\par\vspace{0.7cm}%
	{\large\color{headingcolor}\textbf{#3}}%
	\par\vspace{0.35cm}%
}

\providecommand{\cvletterclosing}[1]{%
	\par\vspace{0.2cm}%
	#1,\par
	\vspace{0.7cm}%
	{\color{headingcolor}\textbf{\brandName}}\par
}

	\input{letters/\cvletter.tex}
\fi

\name{\brandFirstName}{\brandLastName}
\tagline{\cvTagline}%|
% No \photo: dropping it makes \makecvheader give the header block the full
% \linewidth, which is what lets the three columns below span the page. The
% Gauss citation, which used to occupy the photo slot, is now column three.
%
% Contact details stacked in one column ran nine lines deep and set the height
% of the whole header. Split across two columns beside the citation, the block
% is five lines instead -- the saving the one-page layout needs.
%
% Each entry is an \mbox (see \socialtext in the class), so it cannot wrap: a
% column narrower than its longest entry overflows rather than reflowing.
% Widths are therefore sized to the longest line in each column, not evenly.
%
% The class appends \hspace{1em} to every entry to separate entries that share
% a line. Here each entry owns its line, so that 1em is dead width at the right
% edge of every column -- and it was exactly what pushed the longest entries
% past their column. Dropped locally rather than in the class, which other
% documents share.
%
% The icon sits in a fixed-width box rather than being followed by a fixed
% space, so every entry's text starts at the same x no matter how wide its
% symbol is. With a plain \hspace the text began wherever the glyph happened
% to end: the FontAwesome symbols here differ by up to 1.9pt (map marker vs.
% info), so the middle column stepped in and out by that much down its five
% lines. It had been patched per line with hand-tuned \hspace values in the
% entries below -- 0.3em here, 0.5em there, 0.4em on the phone -- which is
% both unmaintainable and was simply missing on the citizenship line, leaving
% it 4.5pt left of its neighbours. Those kerns are gone; do not add new ones.
%
% 1.45em is the widest symbol in this header (the link glyph, 8.33pt at this
% size) plus the 0.5em gap the class used. It costs the widest entry 0.64pt of
% extra width -- the columns are sized to their longest line, so check the
% German build for an overfull box after changing it.
\newcommand{\socialiconwidth}{1.45em}
\renewcommand{\socialtext}[2]{\mbox{\makebox[\socialiconwidth][l]{\textcolor{symbolcolor}{#1}}#2}}
\renewcommand{\sociallink}[3]{\mbox{\makebox[\socialiconwidth][l]{\textcolor{symbolcolor}{#1}}\link{#2}{#3}}}

% Handle instead of full URL for these two: their "linkedin.com/in/" and
% "youtube.com/@" prefixes are already said by the icon beside them, and they
% were the two widest unbreakable entries in the block. Dropping the prefixes
% is what leaves the group narrower than \linewidth -- without it there is no
% spare width and "centred" would be indistinguishable from full-width. The
% links themselves still point at the full URLs.
\renewcommand*{\linkedin}[1]{\sociallink{\linkedinSymbol}{https://www.linkedin.com/in/#1/}{#1}}
\renewcommand*{\youtube}[1]{\sociallink{\youtubeSymbol}{https://youtube.com/@#1}{@#1}}

% \makebox centres the three columns as one group: the gutters are fixed, so
% the leftover width falls outside the group as equal left/right margins
% instead of being pushed between the columns by \hfill.
\socialinfo{
	\makebox[\linewidth][c]{%
	% Sized to the e-mail address, now the longest entry in this column.
	\begin{minipage}[t]{0.26\linewidth}
		\raggedright
		\githubProfile\newline
		\linkedin{\brandLinkedin}\newline
		\brandWebsite\newline
		% YouTube commented out -- no videos published yet, and a channel handle
		% that leads to an empty channel costs more than it pays. Note the
		% \newline moved up to the e-mail line: this entry was last in the
		% minipage and carried no \newline of its own, so commenting it out
		% alone would have left the e-mail's \newline dangling and set an empty
		% final line. Restore both together.
		\email{\brandEmail}
		%\newline
		%\youtube{\brandYoutube}
	\end{minipage}%
	\hspace{1.5em}%
	% Widest column: it carries the German "Geburtstag: ..., Schwäbisch Hall",
	% which is longer than its English counterpart. Sized for the German build.
	\begin{minipage}[t]{0.395\linewidth}
		\raggedright
		\address{Stuttgart}\newline
		\ifcvrelocate
		\infos{\IfLanguageName{english}{Open to relocating to Paris}{Umzug nach Paris möglich}}\newline
		\fi
		\infos{\IfLanguageName{english}{German Citizen}{Deutscher Staatsbürger}}\newline
		% ", Germany" dropped: the citizenship line above already says it, and
		% this is the longest entry in the column.
		% Geburtsort commented out per request -- the date stays, the place is
		% gone. Restore by putting ", Schwäbisch Hall" back inside both branches:
		%   {Birthday: \DTMdisplaydate{1995}{12}{25}, Schwäbisch Hall}
		%   {Geburtstag: \DTMdisplaydate{1995}{12}{25}, Schwäbisch Hall}
		% This was the widest entry in the column and the one the column width
		% was sized for (see the note above the minipage); with it gone the
		% column has slack, so nothing needs re-measuring until it comes back.
		% The empty fourth argument is required, not decoration: \DTMdisplaydate
		% takes {year}{month}{day}{dow}. The old line got away with three brace
		% groups because the ", Schwäbisch Hall" that followed handed it a comma
		% as the day-of-week, which the date style then printed as nothing. Take
		% the place away and the macro swallows the closing brace instead --
		% "Argument of \DTMdisplaydate has an extra }", and no PDF at all.
		% Date of birth commented out on request -- the birthplace went earlier,
		% now the date follows. Costs the layout nothing: this column drops from
		% five entries to four, and the contact column on the left still sets the
		% header height at five. Restore by uncommenting; note the empty fourth
		% argument to \DTMdisplaydate is required (see the note above).
		%\infos{\IfLanguageName{english}{Birthday: \DTMdisplaydate{1995}{12}{25}{}}{Geburtstag: \DTMdisplaydate{1995}{12}{25}{}}}\newline
		\infos{\IfLanguageName{english}{Driving License: Class B}{Führerschein: Klasse B}}\newline
		\smartphone{+49 15253799950}
	\end{minipage}%
	\hspace{1.5em}%
	\begin{minipage}[t]{0.24\linewidth}
		\raggedright
		\footnotesize\itshape
		\IfLanguageName{english}{%
			\textquote[Carl Friedrich Gauss]{Truly, it is not knowledge, but learning; not possession, but acquisition; not being there, but arriving -- that grants the greatest joy.}%
		}{%
			\textquote[Carl Friedrich Gau\ss]{Wahrlich es ist nicht das Wissen, sondern das Lernen, nicht das Besitzen, sondern das Erwerben, nicht das Da-Seyn, sondern das Hinkommen, was den grössten Genuss gewährt.}%
		}%
	\end{minipage}%
	}%
}

\begin{document}

\makecvheader

\makecvfooter
% The left footer field was empty. The most credible line for "I work in the
% open" is the one showing this PDF is itself the artifact of a public
% pipeline -- and in the footer it costs the one-page layout nothing.
{\footnotesize\textsc{Typeset with}\ \link{\brandGithubUrl/DocumANTation}{DocumANTation}}
% "CV" became \cvfootertitle (2026-09-23), so a cover letter can say what it
% is. The line before, kept:
% {\textsc{\brandName{} - CV}}
{\textsc{\brandName{} - \cvfootertitle}}
{\thepage}

% The section list moved into profiles/: which sections fit is a per-target
% decision, and it was already being made here by commenting lines in and out.
\cvBody

\end{document}
